% =============================================================
% reference.tex — 付録B：用語クイックリファレンス
% =============================================================
\chapter*{付録B：用語クイックリファレンス}
\addcontentsline{toc}{chapter}{付録B：用語クイックリファレンス}

試験直前に見返すための用語一覧です。
各分野の頻出用語を1行で整理しています。

% =============================================================
\section*{ストラテジ系の重要用語}
% =============================================================

{
\small
\renewcommand{\arraystretch}{1.3}
\arrayrulecolor{HairLine}
\begin{longtable}{@{}>{\color{Navy}\bfseries\sffamily\footnotesize}p{3.8cm} p{8.2cm}@{}}
\arrayrulecolor{HikakuBar}\toprule
\textbf{用語} & \textbf{一言でいうと} \\
\arrayrulecolor{HairLine}\midrule
\endfirsthead
\arrayrulecolor{HairLine}\midrule
\textbf{用語} & \textbf{一言でいうと} \\
\midrule
\endhead
経営理念 & 企業が存在する根本的な目的・価値観 \\
経営戦略 & 理念を実現するための具体的な道筋 \\
SWOT分析 & 強み・弱み・機会・脅威の4象限で自社を分析 \\
VRIO分析 & 経済価値・希少性・模倣困難性・組織で強みを深堀り \\
PPM & 市場成長率×占有率で事業を4象限に分類 \\
3C分析 & 顧客・競合・自社の3視点で市場を分析 \\
5フォース分析 & 業界の競争圧力を5つの力で分析 \\
4P / 4C & マーケティングミックスの売り手/買い手視点 \\
アンゾフの成長マトリクス & 既存/新規 × 市場/製品で成長方向を4分類 \\
DX & ITでビジネスモデルや企業文化を変革すること \\
PoC & 技術的な実現可能性を確認する小規模な検証 \\
MVP & 市場仮説を検証するための最小限の製品 \\
リーンスタートアップ & MVPで仮説検証を繰り返す手法 \\
損益分岐点（BEP） & 利益がゼロになる売上高 = 固定費 ÷ 限界利益率 \\
売上総利益 & 売上高 $-$ 売上原価（粗利） \\
営業利益 & 売上総利益 $-$ 販管費（本業の儲け） \\
経常利益 & 営業利益 ± 営業外損益 \\
特許権 & 発明を保護（出願から20年） \\
著作権 & 創作的表現を保護（死後70年・無方式主義） \\
偽装請負 & 請負契約で発注者が直接指揮命令する違法状態 \\
\arrayrulecolor{HikakuBar}\bottomrule
\end{longtable}
}

% =============================================================
\section*{マネジメント系の重要用語}
% =============================================================

{
\small
\renewcommand{\arraystretch}{1.3}
\arrayrulecolor{HairLine}
\begin{longtable}{@{}>{\color{Navy}\bfseries\sffamily\footnotesize}p{3.8cm} p{8.2cm}@{}}
\arrayrulecolor{HikakuBar}\toprule
\textbf{用語} & \textbf{一言でいうと} \\
\arrayrulecolor{HairLine}\midrule
\endfirsthead
\arrayrulecolor{HairLine}\midrule
\textbf{用語} & \textbf{一言でいうと} \\
\midrule
\endhead
要件定義 & 「何を作るか」を決める工程（ユーザ主体） \\
外部設計 & 画面・帳票・インターフェースを設計 \\
内部設計 & プログラムの内部構造を設計 \\
ウォーターフォール & 工程を順に進め、後戻りしにくい開発モデル \\
アジャイル開発 & 短い反復で段階的に完成させる開発手法 \\
スクラム & アジャイルの代表的フレームワーク \\
スプリント & スクラムにおける1〜4週間の反復単位 \\
DevOps & 開発と運用の一体化（CI/CD） \\
V字モデル & 設計工程とテスト工程の対応関係 \\
WBS & 作業を階層的に分解した構造 \\
ガントチャート & 横棒で進捗を管理する日程表 \\
クリティカルパス & 最長経路＝プロジェクト全体の最短所要日数 \\
QCD & 品質・コスト・納期の三大制約 \\
SLA & サービス品質の合意文書 \\
SLM & SLAを継続的に管理・改善するプロセス \\
ITIL & ITサービスマネジメントのベストプラクティス集 \\
インシデント管理 & サービスの迅速な復旧（応急処置） \\
問題管理 & インシデントの根本原因を除去し再発防止 \\
MTBF & 平均故障間隔（動いている時間） \\
稼働率 & MTBF ÷ (MTBF + MTTR) \\
システム監査 & 独立した立場でISの信頼性等を評価 \\
職務分掌（SoD） & 業務を分離して不正・ミスを防止 \\
\arrayrulecolor{HikakuBar}\bottomrule
\end{longtable}
}

% =============================================================
\section*{テクノロジ系の重要用語}
% =============================================================

{
\small
\renewcommand{\arraystretch}{1.3}
\arrayrulecolor{HairLine}
\begin{longtable}{@{}>{\color{Navy}\bfseries\sffamily\footnotesize}p{3.8cm} p{8.2cm}@{}}
\arrayrulecolor{HikakuBar}\toprule
\textbf{用語} & \textbf{一言でいうと} \\
\arrayrulecolor{HairLine}\midrule
\endfirsthead
\arrayrulecolor{HairLine}\midrule
\textbf{用語} & \textbf{一言でいうと} \\
\midrule
\endhead
2進数 & 0と1だけで表す数の体系 \\
2の補数 & コンピュータ内部で負の数を表す方式 \\
AND（論理積） & 両方1のときだけ1 \\
OR（論理和） & どちらか一方でも1なら1 \\
XOR（排他的論理和） & 2つが異なるとき1 \\
レジスタ & CPU内部の超高速記憶 \\
キャッシュメモリ & CPUとメインメモリの間の高速バッファ（SRAM） \\
メインメモリ & 実行中プログラムの記憶領域（DRAM） \\
RAID 0 & ストライピング（高速化のみ・冗長性なし） \\
RAID 1 & ミラーリング（同じデータを2台に書込み） \\
RAID 5 & パリティ分散（3台以上・1台故障まで復旧可） \\
OSS & ソースコード公開、自由に利用・改変・再配布可 \\
TCP & コネクション型、信頼性のある通信 \\
UDP & コネクションレス型、高速だが信頼性なし \\
DNS & ドメイン名 ↔ IPアドレスの変換（名前解決） \\
DHCP & IPアドレスの自動割り当て \\
サブネットマスク & ネットワーク部とホスト部の境界を示す \\
主キー & レコードを一意に識別（NULL・重複不可） \\
外部キー & 他テーブルの主キーを参照（テーブル間の関連づけ） \\
SQL（SELECT） & テーブルからデータを取得する命令 \\
正規化 & データの冗長性を排除して更新矛盾を防ぐ \\
CIA三要素 & 機密性・完全性・可用性 \\
公開鍵暗号 & 公開鍵で暗号化、秘密鍵で復号（RSA） \\
デジタル署名 & 送信者の秘密鍵で署名、公開鍵で検証 \\
ランサムウェア & ファイルを暗号化して身代金を要求 \\
フィッシング & 偽サイトで個人情報を騙し取る攻撃 \\
SQLインジェクション & 入力フォームからDB操作する攻撃 \\
フェールセーフ & 故障時に安全な状態で停止する設計 \\
AI / ML / DL & AI $\supset$ ML $\supset$ DL の入れ子構造 \\
AIバイアス & 学習データの偏りによる不公平な判断 \\
\arrayrulecolor{HikakuBar}\bottomrule
\end{longtable}
}

% =============================================================
\section*{直前10パターン}
% =============================================================

\begin{matome}[試験直前に確認する10パターン]
\begin{enumerate}
  \item \textbf{PLの5段階利益}：売上総利益→営業利益→経常利益→税引前当期純利益→当期純利益
  \item \textbf{知的財産権の保護期間}：特許20年、実用新案10年、意匠25年、商標10年（更新可）、著作権 死後70年
  \item \textbf{請負/準委任/派遣}：指揮命令権は請負・準委任＝受注者、派遣＝派遣先
  \item \textbf{V字モデル}：要件定義↔運用テスト、外部設計↔システムテスト、内部設計↔結合テスト、プログラミング↔単体テスト
  \item \textbf{稼働率}：$\text{MTBF} \div (\text{MTBF} + \text{MTTR})$、直列＝積、並列＝$1-(1-A_1)(1-A_2)$
  \item \textbf{TCP vs UDP}：TCP＝信頼性（再送あり）、UDP＝速度（再送なし）
  \item \textbf{正規化}：第1＝繰り返し除去、第2＝部分従属除去、第3＝推移従属除去
  \item \textbf{暗号と署名}：暗号化＝受信者の公開鍵、デジタル署名＝送信者の秘密鍵
  \item \textbf{CIA三要素}：機密性（C）・完全性（I）・可用性（A）
  \item \textbf{AI/ML/DL}：AI $\supset$ ML $\supset$ DL、教師あり/教師なし/強化学習
\end{enumerate}
\end{matome}
